\subsubsection{The general case}
Now let us generalize Proposition \ref{finde1} to the nonsemisimple case. For this purpose assume that every Poisson derivation of $A$ of nonpositive degree is inner. Then the Lie algebra $\g$ of the group
$G:=\operatorname{Aut}_0(A)$ of filtration-preserving Poisson automorphisms of $A$ commuting with $s$ has the form $\g=A_2$, where the operation is the Poisson bracket.

Let us also assume that the group~$G$ is reductive. Then it is easy to see that the action of $\g$ on $A$ uniquely lifts to any quantization ${\mathbf A}$, so we have ${\rm Der}({\mathbf A})^s=\g$. Thus we have a Lie algebra inclusion $\g\subset F_2{\mathbf A}\subset {\mathbf A}$.

Let $g\in G$. Write $g=g_0\exp(e)=\exp(e)g_0$ where $g_0\in G$ is semisimple and $e\in \g$ is nilpotent.

\begin{Proposition}\label{finde2} Under the above assumptions, the conclusion of Proposition~{\rm \ref{finde1}} holds for~$g$. Moreover,
\begin{gather*}
\dim HH_0({\mathbf A},{\mathbf A}g)^s\le \dim HH_0({\mathbf A},{\mathbf A}g_0)^s\le \dim HP_0(X^{g_0})^s<\infty.
\end{gather*}
\end{Proposition}

\begin{proof} Let $Z$ be the centralizer of $g_0$ in $G$, a reductive group containing $\exp(e)$. By the Jacobson--Morozov lemma, the Lie algebra $\mathfrak{z}$ of $Z$ contains an $\mathfrak{sl}_2$-triple
$e$, $h$, $f$. We have a~decomposition ${\mathbf A}=\oplus_{j\in {\mathbb Z}}{\mathbf A}_j$, where ${\mathbf A}_j$ is the $j$-eigenspace of~$h$. Since $g(e)=e$, we have $e{\mathbf b}-{\mathbf b}g(e)=[e,{\mathbf b}]$, so for any $T\in HH_0({\mathbf A},{\mathbf A}g)^*$ and ${\mathbf b}\in {\mathbf A}$ we have $T([e,{\mathbf b}])=0$. By representation theory of $\mathfrak{sl}_2$, this implies that $T({\mathbf a})=0$ for any ${\mathbf a}\in {\mathbf A}_j$ with $j>0$.

Now consider the 1-parameter subgroup $t^h$ of $G$ ($t\in {\mathbb C}^\times$), and define
$T_t({\mathbf a}):=T\big(t^{-h}{\mathbf a}t^h\big)$. If ${\mathbf a}\in {\mathbf A}_j$ then $T_t({\mathbf a})=t^{-j}T({\mathbf a})$. Thus, since $T$ vanishes on ${\mathbf A}_j$ for $j>0$, there exists a limit $T_0=\lim\limits_{t\to 0}T_t$, which coincides with $T$ on ${\mathbf A}_0$ and vanishes on ${\mathbf A}_j$ for $j\ne 0$.

We have $T({\mathbf a})=T_t({\mathbf a}_t)$, where ${\mathbf a}_t:=t^h{\mathbf a} t^{-h}$.
Recall that
\begin{gather*}
{\mathbf a}{\mathbf b}-{\mathbf b} g({\mathbf a})={\mathbf a}{\mathbf b}-{\mathbf b}g_0({\mathbf a})-{\mathbf b}\sum_{i\ge 1}\frac{g_0e^i}{i!}({\mathbf a})\in \operatorname{Ker}(T).
\end{gather*}
Hence
\begin{gather*}
t^h({\mathbf a}{\mathbf b}-{\mathbf b} g({\mathbf a}))t^{-h}={\mathbf a}_t{\mathbf b}_t-{\mathbf b}_t g_0({\mathbf a}_t)-{\mathbf b}_t\sum_{i\ge 1}\frac{t^{2i}e^i}{i!}({\mathbf a}_t)\in \operatorname{Ker}(T_t).
\end{gather*}
Thus, for any ${\mathbf a},{\mathbf b}\in {\mathbf A}$ the element ${\mathbf a}{\mathbf b}- {\mathbf b} g_0({\mathbf a})-{\mathbf b}\sum\limits_{i\ge 1}\frac{t^{2i}g_0e^i}{i!}({\mathbf a})$ belongs to $\operatorname{Ker}(T_t)$. Sending~$t$ to zero, we see that ${\mathbf a}{\mathbf b}-{\mathbf b} g_0({\mathbf a})\in \operatorname{Ker}(T_0)$. Thus, $T_0\in (HH_0({\mathbf A},{\mathbf A}g_0)^s)^*$, which is finite dimensional (with dimension dominated by $\dim HP_0(X^{g_0})^s$) by Proposition~\ref{finde1}.

It remains to show that the assignment $T\mapsto T_0$ is injective; this implies both statements of the theorem. To this end, assume that $T_0=0$ but $T\ne 0$, and let $m>0$ be the order of vanishing of $T_t$ at $t=0$. Let $T_0':=\lim\limits_{t\to 0}T_t/t^m$. Then $T_0'$ is a nonzero element of $HH_0({\mathbf A},{\mathbf A}g_0)^*$ which vanishes on ${\mathbf A}_0$. But if ${\mathbf a}\in {\mathbf A}_j$, $j\ne 0$, then ${\mathbf a}=\frac{1}{j}[h,{\mathbf a}]$, so $T_0'({\mathbf a})=0$. This is a~contradiction, which completes the proof of the proposition.
\end{proof}

\begin{Theorem}\label{finde3} Let $A=\mathcal{O}(X)$, where $X$ is a conical symplectic singularity
with Poisson bracket of degree $-2$ such that the group $G$ of dilation-invariant Poisson automorphisms of $X$ is reductive. Then for any quantization ${\mathbf A}$ of $A$ and any $g\in G$ the space $HH_0({\mathbf A},{\mathbf A}g)^s$ is finite dimensional, with dimension dominated by $\dim HP_0(X^{g_0})^s$, which is finite. Therefore, nondegenerate short star-products on~$A$ depend on finitely many parameters.
\end{Theorem}

\begin{proof}The theorem follows immediately from the above results and Lemmas \ref{inn}, \ref{finfix}.
\end{proof}
